\section{Homogene Koordinaten}

Homogene Koordinaten bieten eine alternative Möglichkeit Punkte und Vektoren im kartesischen Koordinatensystem darzustellen. Im Rahmen des Raytracers spielt dabei vor allem die Darstellung von Punkten
und Vektoren im dreidimensionalen Raum eine wichtige Rolle, weshalb sich nachfolgend alle Erläuterungen auf Elemente des dreidimensionalen Raums beziehen.

In der gängigen Darstellungsweise wird ein Vektor $\mathbf{a}$ durch 3 Komponenten dargestellt
\begin{equation*}
\textbf{a} = \begin{pmatrix}x \\ y\\ z\end{pmatrix}
\end{equation*}

und eine Matrix \textbf{M} besitzt insgesamt 9 Einträge der Form:
\begin{equation*}
\textbf{M} = \begin{pmatrix}
    m_{11} & m_{12} & m_{13}\\
    m_{21} & m_{22} & m_{23}\\
    m_{31} & m_{32} & m_{33}
\end{pmatrix}
\end{equation*}

Auf einen Vektor lässt sich nun eine Reihe von Transformationen anwenden, wie z.B. Rotation oder Translation.
Diese beiden Transformationsarten werden verdeutlichen weshalb man überhaupt homogene Koordinaten verwendet.
Während sich Die Rotation des Vektors $\mathbf{a}$ durch die Multiplikation mit einer Rotationsmatrix $\mathbf{R_x(\phi)}$ (dabei handelt es sich um eine spezielle Rotationsmatrix, die nur um die x-Achse rotiert) bewerkstelligen lässt


\begin{align*}
    \mathbf{R_x(\phi)}& = \begin{pmatrix}
                         1 & 0 & 0 \\
                         0 & \cos \phi & -\sin \phi \\
                         0 & \sin \phi & \cos \phi
                         \end{pmatrix}\\        
    \mathbf{a_{rot}}& = \mathbf{R_x(\phi)} \cdot \mathbf{a}
\end{align*}


muss für die Translation die komponentenweise Addition mit einem Verschiebungsvektor $\mathbf{t}$ durchgeführt werden

\begin{align*}
    \mathbf{a_{trans}} = \mathbf{a} + \mathbf{t} = \begin{pmatrix} a_x + t_x \\ a_y + t_y \\ a_z + t_z
                                                   \end{pmatrix}
\end{align*}

Durch Verwendung von homogenen Koordinaten entfällt diese gesonderte Behandlung unterschiedlicher
Transformationen und es werden zwei neue Objekte eingeführt, nämlich Punkte und Vektoren. In der
bisher beschriebenen Darstellungsform kann nämlich nicht unterschieden werden ob ein kartesischer Vektor nun
zur Repräsentation eines Punktes eines geometrischen Objekts verwendet wird, oder aber als Vektor im eigentlichen
Sinn, das heißt als richtungsangebendes Objekt. Diese Information ist jedoch wichtig, da Punkte verschoben werden, Vektoren(als richtungsangebendes Objekt) allerdings unbeeinflusst bleiben.\\

Nun zur Darstellung in homogenen Koordinaten. Ein homogener Vektor $\mathbf{b}$ besitzt 4 Komponenten.
\begin{equation*}
    \mathbf{b} = \begin{pmatrix} x \\ y \\ z \\ w   \end{pmatrix}
\end{equation*}

 Wie man sieht wurde ein kartesischer Vektor lediglich um die Komponente $w$ erweitert. 
\begin{align*}
    \mathbf{p}& = \begin{pmatrix} x \\ y \\ z \\ 1  \end{pmatrix} \text{entspricht einem Punkt}\\
    \mathbf{v}& = \begin{pmatrix} x \\ y \\ z \\ 0  \end{pmatrix} \text{entspricht einem Vektor}\\
\end{align*}

Diese Komponente 
ermöglicht einerseits die Unterscheidung zwischen Punkt Und Vektor und dient im Falle des Punktes als
Skalierungsfaktor. Das heißt ein Punkt in homogenen Koordinaten lässt sich durch einen kartesischen Vektor
$\mathbf{c}$ der Form
\begin{equation*}
   \mathbf{c} = \begin{pmatrix} x/w \\ y/w \\ z/w   \end{pmatrix}
\end{equation*}
beschreiben.\\

Auch die homogenen Matrizen müssen natürlich um eine Dimension erweitert werden und besitzen nun 16 Einträge.
Sei $\mathbf{N}$ eine homogene Matrix:

\begin{equation*}
    \mathbf{N} = \begin{pmatrix}
        r_{11} & r_{12} & r_{13} & t_{14}\\
        r_{21} & r_{22} & r_{23} & t_{24}\\
        r_{31} & r_{32} & r_{33} & t_{34}\\
        p_{41} & p_{42} & p_{43} & w\\
    \end{pmatrix}
\end{equation*}

Dabei spielen vor allem die Komponenten mit dem Suffix $r$ und $t$ eine wichtige Rolle. Spezielle 
Transformationsmatrizen zum Rotieren, Skalieren, Spiegeln und Scheren benutzen die $r$ Komponenten wohingegen
Translationsmatrizen die $t$ Komponenten verwenden.

Am nachfolgenden Beispiel sollen die beiden angesprochenen Vorteile von homogenen Koordinaten nochmals verdeutlicht werden. Dazu soll der Punkt $\mathbf{p} = \begin{pmatrix} 1 \\ 2 \\ -1 \\ 1 \end{pmatrix}$ und der 
Vektor $\mathbf{v} = \begin{pmatrix} 1 \\ 2 \\ -1 \\ 0 \end{pmatrix}$ um den kartesischen Vektor
$\mathbf{s} = \begin{pmatrix} 3 \\ -1 \\ -1 \end{pmatrix}$ verschoben werden. Dazu wird Die Translationsmatrix
$\mathbf{T}$ erzeugt und mit $\mathbf{p,v}$ multipliziert.

\begin{align*}
    \mathbf{T}& = \begin{pmatrix}
        1 & 0 & 0 & 3\\
        0 & 1 & 0 & -1\\
        0 & 0 & 1 & -1\\
        0 & 0 & 0 & 1
    \end{pmatrix}\\
    \mathbf{p_{trans}}& = \mathbf{T} \cdot \mathbf{p} = \begin{pmatrix} 4 \\ 1 \\ -2 \\ 1  \end{pmatrix}\\
    \mathbf{v_{trans}}& = \mathbf{T} \cdot \mathbf{v} = \begin{pmatrix} 1 \\ 2 \\ -1 \\ 0  \end{pmatrix}
\end{align*}

Zwei Dinge werden deutlich: Erstens lässt sich nun die Translation einheitlich wie die anderen Transformationen
durch eine Matrixmultiplikation durchführen. Und zweitens bleibt, wie bereits erwähnt, der Vektor als rein 
richtungsangebendes Objekt von der Translation verschont.\\


Weiterhin sind auch noch Addition und Subtraktion zwischen homogenen Vektoren erlaubt jedoch nicht in beliebiger
Kombination, da die die Addition zweier Punkte beispielsweise zu keinem vernünftig interpretierbaren Ergebnis führen würde. Dies ist aus Tabelle~\ref{tab_Operationen} zu entnehmen

\begin{table}[ht]
\centering
\begin{tabular}{|c c|}
\hline 
Operation & Ergebnis \\ 
\hline 
Vektor + Vektor  & Vektor \\ 
\hline 
Vektor - Vektor & Vektor \\ 
\hline 
Punkt - Punkt & Vektor \\ 
\hline 
Punkt + Vektor & Punkt \\ 
\hline 
Punkt - Vektor & Punkt \\ 
\hline 
\end{tabular} 
\caption{Erlaubte Additions- und Subtraktionsoperationen, sowie das entsprechende Ergebnis}
    \label{tab_Operationen}
\end{table}